\documentclass[11pt]{article}
\usepackage[margin=0.85in]{geometry}
\usepackage[T1]{fontenc}
\usepackage{lmodern}
\usepackage{microtype}
\usepackage{booktabs}
\usepackage{tabularx}
\usepackage{xcolor}
\usepackage{hyperref}
\usepackage{siunitx}
\usepackage{enumitem}

\definecolor{accent}{HTML}{6A3FC7}
\definecolor{ink}{HTML}{202124}
\hypersetup{colorlinks=true,linkcolor=accent,urlcolor=accent}
\color{ink}
\sisetup{detect-all,group-separator={,},group-minimum-digits=4}
\setlist{nosep,leftmargin=*}
\setlength{\parskip}{0.55em}
\setlength{\parindent}{0pt}

\title{\textbf{wgfx Performance Relative to Raw WebGPU}\\
\large Order-Controlled Microbenchmark Evidence}
\author{e10b / wgfx project}
\date{10 October 2026}

\begin{document}
\maketitle

\begin{abstract}
This note evaluates whether the \texttt{wgfx} C++ convenience layer adds
material CPU overhead relative to direct WebGPU calls through
\texttt{wgpu-native}. In order-controlled microbenchmarks, the measured
differences for compute encoding and submission, storage updates, bind-group
creation, and render encoding are small, inconsistent in sign, or within the
observed run-to-run variability. For the operations tested, no repeatable
abstraction penalty is resolved: \texttt{wgfx} is effectively at performance
parity with raw WebGPU.
\end{abstract}

\section{Claim}

The defensible conclusion is:

\begin{quote}
\textbf{For the operations and workloads measured, wgfx is effectively as fast
as raw WebGPU through wgpu-native. No repeatable core forwarding or encoding
overhead was resolved above the benchmark's noise floor.}
\end{quote}

This is not a claim that every possible program, resource-loading pattern, or
backend must behave identically. It is a measured statement about the tested
hot-path operations, under matched resources and workloads.

\section{Experimental setup}

\begin{itemize}
  \item Apple M3 Pro, 14 GPU cores, arm64; macOS 15.6.1.
  \item AppleClang 17.0.0; \texttt{RelWithDebInfo} build
        (\texttt{-O2 -g -DNDEBUG}).
  \item \texttt{wgpu-native} supplied by WebGPU-distribution commit
        \texttt{50010b083e4a1bbdac8f9102602df915fe7404cb}.
  \item Raw and wgfx paths share the same process, device, queue, shaders,
        buffers, textures, pipelines, bind-group shapes, and workload sizes.
  \item Each test performs ten warm-up batches followed by repeated timed
        samples using \texttt{std::chrono::steady\_clock}.
  \item Submission and render tests reverse raw/wgfx execution order on every
        sample. This prevents thermal, cache, and queue drift from consistently
        favoring either implementation.
  \item Five complete suite executions were collected. Values below average
        the per-run medians rather than treating all samples as independent.
\end{itemize}

Queue completion is performed outside CPU submission samples so neither path
inherits greater queue backpressure. The raw and wrapped paths ultimately issue
the same WebGPU operations; the comparison isolates wrapper-side work.

\section{Order-controlled results}

\begin{table}[ht]
\centering
\caption{CPU time per operation. Small positive and negative differences are
measurement variability, not evidence that either path intrinsically
accelerates WebGPU.}
\begin{tabularx}{\textwidth}{@{}X r r r@{}}
\toprule
Operation & Raw WebGPU & wgfx & Relative difference \\
\midrule
Compute encoding, one dispatch & \SI{6.142}{\micro\second} &
  \SI{6.138}{\micro\second} & $-0.1\%$ \\
Storage update, \SI{4}{\kibi\byte} & \SI{2.856}{\micro\second} &
  \SI{2.921}{\micro\second} & $+2.3\%$ \\
Bind-group recreation & \SI{1.766}{\micro\second} &
  \SI{1.717}{\micro\second} & $-2.8\%$ \\
Compute-pipeline creation & \SI{96.925}{\micro\second} &
  \SI{93.584}{\micro\second} & $-3.4\%$ \\
Encode + queue submit, one dispatch & \SI{21.321}{\micro\second} &
  \SI{21.869}{\micro\second} & $+2.6\%$ \\
Submit + completion, one dispatch & \SI{1.197}{\milli\second} &
  \SI{1.215}{\milli\second} & $+1.5\%$ \\
Render encoding, one draw & \SI{17.667}{\micro\second} &
  \SI{16.708}{\micro\second} & $-5.4\%$ \\
Render encoding, 100 draws & \SI{33.717}{\micro\second} &
  \SI{34.108}{\micro\second} & $+1.2\%$ \\
Render encoding, 1,000 draws & \SI{154.442}{\micro\second} &
  \SI{158.950}{\micro\second} & $+2.9\%$ \\
\bottomrule
\end{tabularx}
\end{table}

The most suspicious original result was encode plus submission at
$+22.4\%$ for wgfx. That benchmark always ran raw first and wgfx second. Against
canonical wgfx, after reversing order on every sample and increasing the sample
count, the result became $+2.6\%$. Compute encoding alone measures $-0.1\%$,
submit plus completion measures $+1.5\%$, and 1,000-draw rendering measures
$+2.9\%$. These values resolve no material steady-state wrapper overhead.

Negative values are not interpreted as wgfx being faster. Creation is relatively
noisy and is not a steady-state frame-loop operation. Profiling also found and
removed unconditional diagnostic output from resource creation and repaired
replaced bind-group and pipeline-layout handle ownership before the final runs.
A separate 256 KiB storage create/destroy case retains a 22.3\% wrapper cost;
that resource-lifecycle result is disclosed in the CSV but is distinct from the
steady-state forwarding and encoding claim.

\section{Application exercise: wDoom}

The microbenchmarks are accompanied by a real consumer, wDoom, which uses wgfx
to drive a WebGPU compute implementation of DOOM wall-column rasterization.
DOOM's CPU performs BSP traversal and emits visible wall commands; the compute
shader consumes compact per-column command ranges and indexed texture and
colormap resources. This exercises persistent storage buffers, repeated buffer
updates, bind groups, compute dispatch, queue submission, and presentation in a
complete application rather than an isolated loop.

Correctness was checked with a deterministic software-renderer reference and
GPU buffer readback across all four embedded Ultimate DOOM demos. The run
completed 8,738 eligible frames and compared 559,232,000 indexed pixels with
zero pixel mismatches, zero readback failures, and zero callbacks pending at
completion. This does not independently measure wrapper overhead, but it shows
that the benchmarked wgfx operations are used successfully in a substantial,
pixel-validated rendering workload.

\section{Limitations}

\begin{itemize}
  \item Results apply to this machine, backend revision, compiler, and workload.
  \item CPU wall-clock timing cannot isolate shader-only GPU duration; GPU
        timestamp queries were unavailable in the requested feature set.
  \item Texture decoding, mip generation, unusual allocation patterns, and
        other APIs were not proven equivalent by these hot-path measurements.
  \item Differences of a few percent that change sign or move with test order
        should be treated as unresolved, not promoted into performance claims.
\end{itemize}

\section{Conclusion}

The evidence supports saying that \texttt{wgfx} is about as fast as direct
WebGPU for the core operations tested. Once test-order bias was removed, no
repeatable encoding or submission overhead remained. The order-controlled
results cluster within a few percent of raw WebGPU and change sign across
operations and runs. At this measurement resolution, those differences are
noise rather than evidence of a systematic wrapper cost.

\section*{Reproducibility artifacts}

The repository contains the benchmark source, full methodology, and CSV data:

\begin{itemize}
  \item \texttt{benchmarks/wgfx\_benchmark.cpp}
  \item \texttt{PERFORMANCE.md}
  \item \texttt{benchmarks/results/m3-pro-order-controlled-summary.csv}
  \item wDoom's \texttt{wgpu} branch and four-demo validation corpus
\end{itemize}

\end{document}
